% ---------------------------------------------------------------------------
% Appendix A: Supplementary Experimental Results
% ---------------------------------------------------------------------------
\section{Supplementary Experimental Results}
\label{sec:appendix-results}

In this appendix we use the same terminology as in the main paper: \emph{page-level} retrieval indexes each page separately and aggregates scores at query time. \emph{document-level} retrieval composes a grid per document and encodes it once. We distinguish \textbf{First-4 pages} (first four pages used in page-level aggregation) from \textbf{First-4 grid} (first four pages composed into a single document-level grid). All NDCG@5 values are reported in percentage (0 to 100).


% ---------------------------------------------------------------------------
\subsection{Cross-Strategy Generalization of DocPC}
\label{sec:app-a2}

Table~\ref{tab:app-a3} reports cross-strategy generalization of DocPC-ColQwen. Rows denote training page-selection strategies and columns denote test grid strategies. Values are average NDCG@5 (\%) over the seven domains.

\begin{table}[htbp]
  \centering
  \footnotesize
  \setlength{\tabcolsep}{4pt}
  \caption{Cross-strategy generalization of DocPC-ColQwen (average NDCG@5, \%). Rows denote training strategies and columns denote test strategies. F4 = First-4, Bnd. = Boundary, U4 = Uniform-4, R4 = Random-4, and L4 = Last-4.}
  \label{tab:app-a3}
  \begin{tabular}{@{}lccccc@{}}
    \toprule
    Train $\backslash$ Test & F4 & Bnd. & U4 & R4 & L4 \\
    \midrule
    Bnd. & 42.78 & 40.75 & 40.34 & 38.62 & 36.33 \\
    F4   & 40.66 & 39.17 & 38.95 & 36.77 & 35.36 \\
    L4   & 40.23 & 38.28 & 38.72 & 37.00 & 34.93 \\
    R4   & 42.64 & 40.69 & 40.17 & 39.60 & 36.00 \\
    U4   & \textbf{43.70} & 41.33 & 41.29 & 38.56 & 36.90 \\
    \bottomrule
  \end{tabular}
\end{table}

% ---------------------------------------------------------------------------
% ---------------------------------------------------------------------------
\subsection{Comparison of Tuned ColQwen and DocPC Under Different Test-Time Strategies}
\label{sec:app-a2b}

Table~\ref{tab:app-a3b} compares tuned ColQwen and DocPC-ColQwen under different test-time representative-page strategies. All values are average NDCG@5 (\%) over the seven domains.

\begin{table}[t]
  \centering
  \footnotesize
  \setlength{\tabcolsep}{4pt}
  \caption{Tuned ColQwen vs.\ DocPC-ColQwen under different test-time representative-page strategies on DocViRe (average NDCG@5, \%).}
  \label{tab:app-a3b}
  \begin{tabular}{@{}llc@{}}
    \toprule
    Method & Strategy & Avg \\
    \midrule
    ColQwen (DocViRe-sft) & First-4   & 40.75 \\
    ColQwen (DocViRe-sft) & Boundary  & 39.26 \\
    ColQwen (DocViRe-sft) & Uniform-4 & 37.96 \\
    ColQwen (DocViRe-sft) & Random-4  & 36.23 \\
    \midrule
    DocPC-ColQwen         & First-4   & \textbf{43.70} \\
    DocPC-ColQwen         & Boundary  & 41.33 \\
    DocPC-ColQwen         & Uniform-4 & 41.29 \\
    DocPC-ColQwen         & Random-4  & 39.60 \\
    \bottomrule
  \end{tabular}
\end{table}

% ---------------------------------------------------------------------------
\subsection{Extended Ablations on Training Objective and Page Budget}
\label{sec:app-a4}

The following tables are \emph{extended/exploratory} ablations. Configurations and training setup may differ from the main paper. Values are not directly comparable to Table~\ref{tab:main-results} or the page-budget table in Section~\ref{sec:ablation}, which use controlled settings. The results suggest that the benefit of the listwise term depends strongly on how frequently it is applied. A possible explanation is that the multi-positive contrastive loss primarily shapes the representation space, whereas the listwise ApproxNDCG term mainly refines the ordering among already competitive candidates. When applied too frequently, the listwise term may interfere with early-stage representation learning; when applied sparsely, it can serve as a lightweight ranking refinement on top of a stronger representation space.

\begin{table}[htbp]
  \centering
  \footnotesize
  \caption{Extended loss-frequency ablation for DocPC-ColQwen (NDCG@5, \% AVG). Exploratory. Not directly comparable to main-text Table~\ref{tab:loss-ablation} (controlled configuration). Listwise only / Multi-positive only and Combined($m$) with various $m$.}
  \label{tab:app-a6}
  \begin{tabular}{@{}lc@{}}
    \toprule
    Variant & Avg \\
    \midrule
    Combined (m=1) & 40.53 \\
    Combined (m=3) & 41.61 \\
    Combined (m=5) & 40.66 \\
    Combined (m=7) & 42.70 \\
    Combined (m=9) & 42.96 \\
    Combined (m=11) & 43.18 \\
    Combined (m=13) & 43.67 \\
    Combined (m=15) & 43.83 \\
    Combined (m=20) & 43.36 \\
    Combined (m=30) & 43.27 \\
    Combined (m=40) & \textbf{44.09} \\
    Combined (m=50) & 42.83 \\
    Listwise only & 37.32 \\
    Multi-positive only & 42.68 \\
    \bottomrule
  \end{tabular}
\end{table}

\begin{table}[htbp]
  \centering
  \footnotesize
  \caption{Extended representative-page budget ablation for DocPC-ColQwen (NDCG@5, \% AVG). Exploratory. Main-text page-budget table (Section~\ref{sec:ablation}) uses a different controlled setup (First-1/First-4/First-16, CLIP-selected 4-page grid).}
  \label{tab:app-a7}
  \begin{tabular}{@{}lc@{}}
    \toprule
    Model & Avg \\
    \midrule
    DocPC-ColQwen (First-1) & 39.03 \\
    DocPC-ColQwen (First-4) & 40.66 \\
    DocPC-ColQwen (First-16) & 36.59 \\
    DocPC-ColQwen (CLIP-selected 4-page grid) & 40.73 \\
    \bottomrule
  \end{tabular}
\end{table}

% ---------------------------------------------------------------------------

\subsection{Detailed efficiency breakdown}
\label{sec:app-a4-efficiency}

Table~\ref{tab:app-efficiency-footprint} reports the detailed efficiency and storage comparison between DocPC and page-level ColQwen on the DocViRe test index.
DocPC reduces the number of indexed images and vectors by about $10.1\times$, lowers Milvus storage from 22.4~GB to 2.2~GB, and reduces end-to-end indexing time from 15{,}588~s to 2{,}034~s.
While the composed 4-page grid slightly increases per-image encoding time, the overall encoding and insertion costs are substantially reduced because far fewer inputs need to be indexed.

\begin{table}[htbp]
  \centering
  \footnotesize
  \setlength{\tabcolsep}{4pt}
  \caption{Detailed efficiency and storage comparison between DocPC and page-level ColQwen on the DocViRe test index.}
  \label{tab:app-efficiency-footprint}
  \begin{tabular}{@{}lccc@{}}
    \toprule
    Metric & DocPC & ColQwen & Ratio \\
    \midrule
    Images   & 3{,}850   & 38{,}746 & $10.1\times$ fewer \\
    Vectors  & 2.91M     & 29.25M   & $10.1\times$ fewer \\
    Enc./img & 431.7 ms  & 304.7 ms & 1.42$\times$ higher \\
    Enc.     & 1{,}663 s & 11{,}805 s & $7.1\times$ faster \\
    Insert   & 371 s     & 3{,}783 s & $10.2\times$ faster \\
    Total    & 2{,}034 s & 15{,}588 s & $7.7\times$ faster \\
    Storage  & 2.2 GB    & 22.4 GB  & $10.1\times$ smaller \\
    \bottomrule
  \end{tabular}
\end{table}
% ---------------------------------------------------------------------------

\subsection{Resolution sensitivity}
\label{sec:app-a4-resolution}

We further examine how DocPC-ColQwen (First-4 grid) behaves under different input resolutions of the composed grid image on DocViRe. As shown in Table~\ref{tab:app-a7b}, average NDCG@5 rises from 6.68 at $128{\times}166$ to 22.13 at $256{\times}331$, 37.55 at $512{\times}662$, and 40.65 at the original-resolution setting, indicating consistent but diminishing gains as resolution increases.

\begin{table}[htbp]
  \centering
  \footnotesize
  \caption{Resolution sensitivity of DocPC-ColQwen (First-4 grid) on DocViRe (average NDCG@5, \%).}
  \label{tab:app-a7b}
  \begin{tabular}{@{}lc@{}}
    \toprule
    Resolution & Avg NDCG@5 \\
    \midrule
    $128{\times}166$ & 6.68 \\
    $256{\times}331$ & 22.13 \\
    $512{\times}662$ & 37.55 \\
    Original         & 40.65 \\
    \bottomrule
  \end{tabular}
\end{table}

% ---------------------------------------------------------------------------
\subsection{Retrieval Performance Under Varying Candidate-Pool Sizes}
\label{sec:app-a5}

Table~\ref{tab:app-a8} reports retrieval performance as the candidate document pool size varies (per-domain pool size: 550, 400, or 200 size). The page-level baseline is ColQwen-page with MaxP aggregation (taking the better of the two ColQwen-page variants for each pool size). DocPC-ColQwen (First-4) uses the First-4 grid. ``Gain'' is the absolute difference in average NDCG@5 (percentage points) between DocPC and the page-level baseline.

\begin{table}[htbp]
  \centering
  \footnotesize
  \caption{Retrieval performance under varying candidate-pool sizes (NDCG@5, \% AVG). Page-level baseline: ColQwen-page. Gain = DocPC minus page-level (absolute percentage points).}
  \label{tab:app-a8}
  \begin{tabular}{@{}lccc@{}}
    \toprule
    Pool size & Page-level & DocPC-ColQwen (First-4) & Gain \\
    \midrule
    550 & 40.75 & 44.09 & 3.34 \\
    400 & 41.19 & 44.28 & 3.09 \\
    200 & 50.15 & 52.29 & 2.14 \\
    \bottomrule
  \end{tabular}
\end{table}

% ---------------------------------------------------------------------------
\subsection{Baseline Construction and Scoring Details}
\label{sec:app-baselines}

For text-only baselines, we first extract OCR text from each document. \textsc{BM25} is applied directly to the extracted text, while \textsc{BGE-M3} encodes OCR-derived text chunks into dense embeddings. We also include LLM-summary variants where each document is first compressed into a single abstractive summary and retrieval is performed over summaries using \textsc{BGE-M3} or \textsc{BM25}. OCR outputs are split into fixed-size chunks, each chunk is encoded once, and document scores are obtained by taking the maximum similarity over all chunks (MaxP aggregation).

For page-level visual baselines, each page is treated as a retrieval unit. We compute query-page similarities with late interaction and obtain document scores via MaxP over pages, ensuring a uniform aggregation scheme across page-level methods. We report zero-shot ColQwen-page and ColQwen-page (Boundary), as well as DocViRe-tuned page-level variants: \emph{ColQwen-page (DocViRe-sft, All pages)} and \emph{ColQwen-page (DocViRe-sft, Boundary)}, which use the same supervision as DocPC while retaining page-level indexing and MaxP aggregation during inference.

For document-level bi-encoder baselines, \textsc{Jina-CLIP} and \textsc{Nomic-Embed} are evaluated with First-4 grids, so that each document is represented by a single composed image.

% ---------------------------------------------------------------------------
% Appendix B: Supplementary Dataset Statistics
% ---------------------------------------------------------------------------
\section{Supplementary Dataset Statistics}
\label{sec:appendix-dataset}

\subsection{Dataset Collection and Preprocessing}
To construct DocViRe, we used the first 40 source files to build the training split and files 41 to 50 to build the test split. After obtaining the raw PDFs, we first retained only documents with 4 to 30 pages. We then applied DeepSeek-based document classification and kept only those PDFs assigned to one of the seven target domains: biology, education, finance, government, industrial, legal, and research. The same preprocessing and filtering procedure was applied independently to the test split. This classification step was used only for corpus filtering, rather than as benchmark relevance annotation.

\subsection{Quality validation agreement}
For the test split, the three annotators labeled 98.8\%, 99.0\%, and 94.3\% of checked query-document pairs as relevant, with a three-way full agreement of 92.5\% and an at-least-two agreement of 98.8\%. For query validity, they exhaustively reviewed all test-split queries and marked 99.3\%, 99.7\%, and 100.0\% of them as valid. Across all test queries that received validity labels from all three annotators, the three-way agreement reached 99.0\%.